\begin{claim}{C.52}{Proved}{Cocompact frontier extensions relative to the exponential region}
Use the actual based ambient-null plane loop $f_t$, the groups
$G=\langle f_t\rangle$ and $H$, and the closed set $X_*$ with its
contraction $E_{s,t}$ from C.50.
For every nonopen $G$-orbit $B$, define a based family
$R^B_{s,t}$ as follows:
\[
 R^B_{s,t}=
 \begin{cases}
 \id,&B\text{ a point},\\
 \text{the exponential fraction with parameter }1-s,
       &G|_B\text{ metabelian},\\
 \text{the C.27--C.29 median family},
       &G|_B\text{ nonmetabelian}.
 \end{cases}                                             \tag{C.52.1}
\]
The exponential fraction uses the actual intrinsic Lie image
classified in C.37. On a circle it uses the entire based lift
of the original loop to the covering group, including its
real angle at each intermediate time.

For every open orbit $O$ disjoint from $X_*$ whose full family
$\mathscr F_O$ of actual lifted boundary partitions has
cardinality at most two, there is an intrinsic based
null-homotopy $K^O_{s,t}$ commuting with its original operator
stabilizer and satisfying every actual boundary quotient
equation
\[
                 \Phi_zK^O_{s,t}=R^{Gz}_{s,t}\Phi_z.      \tag{C.52.2}
\]
For each such $O$ with compact proper-operator quotient
$O/Q_O$, this family extends to a based contraction through
homeomorphisms of $\overline O$, jointly continuous with
its inverse, and equal to (C.52.1) on its entire frontier.
It agrees with $E_{s,t}$ on $\overline O\cap X_*$.

These extensions can be chosen with exact transport under
$H$. They agree on common frontiers. Consequently they
paste with the prescribed C.50 contraction on $X_*$ over
every finite union
\[
                  X_*\ \cup\!
                  \bigcup_{O\in\mathcal I}\overline O,    \tag{C.52.3}
\]
where $\mathcal I$ is a finite collection of the
cocompact orbits just specified.
No continuity of (C.52.1) on the entire nonopen locus,
and no extension over an infinite collection or over the
other noncocompact open orbits, is asserted.
\end{claim}

\begin{proof}
\emph{The precise intrinsic target family.}
For a nontrivial intrinsic line or circle orbit, $G|_B$ is
path connected and transitive. The proof of (C.37.9)
classifies every such metabelian image, with its actual
intrinsic homeomorphism-group topology: on a line it is
conjugate to all translations or to the full positive
affine group, and on a circle to all rotations.
That proof uses only these three stated properties of
the image, so it applies to every metabelian orbit in
(C.52.1), without requiring that orbit to lie in $X_*$.

For a real lift $F:X\to\R$ of an actual boundary quotient,
let $b_t$ be its based target path, characterized by
$Fa_t=b_tF$ for the based source lift $a_t$.
The exact source and target coverings are those of
C.42--C.44. The path $b_t$ is a based loop: for a circle
target this is the zero-winding conclusion of C.29,
and for a line target it follows from the identity
endpoint of $f_t$.
Its generated group
\[
                         T=\langle b_t:0\leq t\leq1\rangle
                                                               \tag{C.52.4}
\]
is transitive on $\R$. Indeed its orbits are intervals,
as each word is joined to the identity by a path of
words. A finite endpoint of such an orbit interval is
a common fixed point, which would project to a common
fixed point of the nontrivial transitive target action.
Thus no proper orbit interval is possible.

For a line target, (C.52.4) is the conjugate of the
actual image $G|_B$. For a circle, projection gives a
surjection $T\to G|_B$. When the latter image is
metabelian, its rotation conjugacy lifts to a real
coordinate in which every $b_t$ is a translation.
Transitivity then makes $T$ the full translation group.
Conversely a metabelian $T$ has metabelian quotient
$G|_B$. This proves the exact equivalence used in the
case selection below.

Define, in the specified real lift,
\[
 B_{s,t}=
 \begin{cases}
 \exp_T\bigl((1-s)\log_T b_t\bigr),&T\text{ metabelian},\\
 \text{the C.27 median of the path }b_t,&T\text{ nonmetabelian}.
 \end{cases}                                             \tag{C.52.5}
\]
In the first case $T$ is a simply connected translation
or positive affine Lie group with a global exponential
and logarithm. Formula (C.52.5) and its inverse are
jointly continuous, with all four required based-edge
identities. They commute with every target homeomorphism
commuting with the original path: in the exponential
case each new value lies in $T$, and in the median case
this is C.27.

This rule is natural under any line reparametrization,
including a decreasing one. Such conjugation identifies
the actual groups (C.52.4) topologically, preserves
metabelianness, and, in the Lie cases, intertwines
exponential maps and their unique logarithms. The
median case has the same naturality by C.27.
The source of a real lift of a circle may also be
changed by a covering conjugacy; the same argument
retains the specified based target path.
In particular the resulting map on $B$ is exactly
(C.52.1), independent of all these coordinates.

For clarity, at a time when $f_t|_B=\id$ on a circle,
its based real lift can still be a nonzero integer
translation. Formula (C.52.5) scales that real
translation. It does not replace it by zero.
The endpoint $t=1$ is zero by C.29. This is the same
covering convention as C.50.

\emph{Replacement in the two-foliation construction.}
For $|\mathscr F_O|=2$, retain exactly the source cover,
the two real quotient coordinates, and the domain $D$
of C.45--C.46. Replace each of the two coordinate
median families in (C.46.4) by (C.52.5).
On $D=\R^2$ this is a product of increasing line
homeomorphisms, with a jointly continuous inverse.
On $D=\{u<v\}$ the original two coordinate loops are
identical. Their generated groups and their selected
families are therefore identical, so the strict
inequality is preserved. On $D=\{u<v<u+1\}$ the same
argument applies, and the selected maps commute with
unit translation. In the exponential case this holds
because that translation centralizes $T$; in the
median case it is the original median naturality.
Thus both inequalities are preserved, as are they
by the inverse family.

Every original stabilizing operator, and every one
of its source lifts, permutes the full two-element
family of partitions by the exact comparison
(C.46.5). In these coordinates it is a pair of line
homeomorphisms, possibly followed by exchange.
Commutation with the based source loop identifies
the corresponding target paths by these same
conjugacies. Rule (C.52.5) is natural under them
and makes the same choice on conjugate paths.
Consequently the modified source family commutes
with every such lift, including all source deck maps.
It descends to $O$ with continuous inverse.

Every other actual real quotient is a line
homeomorphism composed with one of the two selected
coordinates, by C.45. The preceding naturality
therefore proves its exact lifted equation too.
Descending through $q_zF_z=\Phi_zp$ proves (C.52.2),
retaining the original source and target covers.

\emph{Replacement in the one-foliation construction.}
For $|\mathscr F_O|=1$, keep the actual closed proper
operator group, its group of all source lifts, and
every coordinate change of C.48.
In (C.48.19) replace only the median base map by
(C.52.5), leaving the fiber interpolation as displayed:
\[
 (u,r)\longmapsto
       \bigl((1-s)h_t(u,r)+su,\ B_{s,t}(r)\bigr).        \tag{C.52.6}
\]
Here the choice of target coordinate is the exact
reparametrization already established in C.48.
We check why all of its operator equations survive
this replacement.

For the unit fiber translation, the fiber equation
$h_t(u+1,r)=h_t(u,r)+1$ is unchanged by the base
replacement. For an operator $(\epsilon u,r+d)$,
the equations are
$h_t(\epsilon u,r+d)=\epsilon h_t(u,r)$ and
$b_t(r+d)=b_t(r)+d$. The first is preserved by
the same convex interpolation; the second says
that translation by $d$ centralizes $T$, so it
holds for (C.52.5).
For the real fiber kernel and a further
$(\lambda u,r+d)$, the fiber function is
$h_t(u,r)=u+\phi_t(r)$ with
$\phi_t(r+d)=\lambda\phi_t(r)$. Scaling $\phi_t$
preserves this equality for either sign of $\lambda$,
and the base comparison is again the same one.

In the normal forms with all real base translations
or a dense lattice base image, the original base
loop is $b_t(r)=r+c_t$. Its selected family is
$B_{s,t}(r)=r+(1-s)c_t$ and commutes with every base
operator. The fiber equations, respectively
independence of $r$ or the exact affine lattice
relation (C.48.21), are preserved by (C.52.6).
The trivial operator case imposes no extra equation.
These are the C.48 normal forms outside its
transitive two-dimensional case. That latter case
already lies in $U_2\subset X_*$ and is absent from
the open orbits of the present claim.

Thus the modified family commutes with the exact
group of all operator lifts, with its original
restriction kernel retained downstairs.
The fiber and base inverse-continuity argument
following (C.48.22) applies unchanged to (C.52.6):
the base family has continuous inverses by
(C.52.5), and fixed bracketing fiber values bound
the increasing convex-interpolation inverses.
The homotopy descends to $O$, with the asserted
endpoint identities. All other actual quotients
are reparametrizations of the one selected
partition; naturality proves (C.52.2) for them.

For an empty family, C.32--C.33 supply a based
null-homotopy in the exact proper-operator
centralizer. There is then no nontrivial actual
quotient equation to check. This finishes the
intrinsic assertion for all three cardinalities.

\emph{Cocompactness and the entire frontier.}
Now assume $O/Q_O$ compact and put $Q=Q_O$.
The first part of C.51 supplies a compact
$C\subset O$ with $QC=O$. It also proves that
every nontrivial frontier orbit gives an actual
quotient in the specified family, after the
explicit basepoint comparison (C.51.8).
In particular an empty family makes all frontier
points fixed.

Define $\overline K^O$ by the just-constructed
$K^O$ on $O$ and by (C.52.1) on $\partial O$.
We first prove its limit along points of $O$;
we do not assume transverse continuity of
(C.52.1). Let
\[
 x_n\in O,\quad x_n\to z\in\partial O,\qquad
                       (s_n,t_n)\to(s,t).              \tag{C.52.7}
\]
Write $x_n=q_ny_n$ with $y_n\in C$ and pass to
a subsequence $y_n\to y\in C$.
The compact set consisting of $C$ and all its
images and inverse images under $K^O$ is exactly
the type of set used in (C.51.11).
Approximate $q_n$ uniformly on this set to
accuracy $1/n$ by restrictions of actual
$h_n\in H_O$, as in (C.51.12).
The open evaluation map of $L=\overline G$ gives
$l_n\to1$ with $l_ny=y_n$. Exact commutation
gives $h_ny=l_n^{-1}h_ny_n\to z$.

The C.42 evaluation homeomorphism, with the
basepoint comparison (C.51.8), now gives
$h_n|_O\to\Phi_z^y$ uniformly on compact subsets.
For a nontrivial target this is an actual
quotient covered by (C.52.2); for a fixed target
it is constant. Commutation with $Q$ therefore
gives precisely
\[
 \begin{split}
 K^O_{s_n,t_n}(x_n)
   &=q_nK^O_{s_n,t_n}(y_n)\\
   &\longrightarrow
       \Phi_z^y K^O_{s,t}(y)
        =R^{Gz}_{s,t}(z). 
 \end{split}                                             \tag{C.52.8}
\]
The approximation error in the second line is
less than $1/n$ before passage to the limit.
Compactness of $C$ and uniqueness of the value
give (C.52.8) for the entire sequence.
This repeats C.51's exact operator comparison
with the new quotient equation; no element of
the restricted closure is assigned an
unproved global extension.

\emph{The diagonal argument at frontier points.}
Take now $z_n\in\partial O$ with $z_n\to z$
and $(s_n,t_n)\to(s,t)$.
For each fixed triple $(s_n,t_n,z_n)$,
(C.52.8) lets us choose $x_n\in O$ such that
\[
 |x_n-z_n|<1/n,\qquad
 \bigl|K^O_{s_n,t_n}(x_n)
            -R^{Gz_n}_{s_n,t_n}(z_n)\bigr|<1/n.         \tag{C.52.9}
\]
Then $x_n\to z$, so applying (C.52.8) to this
new sequence proves
$R^{Gz_n}_{s_n,t_n}(z_n)\to R^{Gz}_{s,t}(z)$.
This establishes the needed joint continuity
on this entire frontier, even across different
target types in (C.52.1).
Splitting a sequence into its interior and
frontier terms proves joint continuity on
$[0,1]^2\times\overline O$.

The same proof applies to the inverse family.
The compact approximation set already contains
all inverse images, and (C.52.2) implies the
exact inverse equation
$\Phi_z(K^O_{s,t})^{-1}=(R^{Gz}_{s,t})^{-1}\Phi_z$.
Use this equation in (C.52.8), then repeat the
diagonal choice (C.52.9) for inverses.
Thus the piecewise inverse is jointly continuous.
All maps are bijections on $O$ and on each
frontier orbit, with these actual inverses;
they are homeomorphisms of the closure.

\emph{Agreement with the fixed exponential region.}
Every group image on an open orbit of $U_2\cup U_3$
is metabelian, by C.36 and C.49. Hence, for all
$g_1,g_2,g_3,g_4\in G$, the plane homeomorphism
$[[g_1,g_2],[g_3,g_4]]$ fixes $U_2\cup U_3$
pointwise. Continuity makes it fix $X_*$ pointwise.
The set $X_*$ is $G$-invariant. Therefore every
nonopen orbit $B\subset X_*$ has metabelian
actual image $G|_B$.

An open orbit outside $U_2\cup U_3$ is disjoint
from $X_*$, because it is an open neighborhood
disjoint from that union. Thus the intersections
$\overline O\cap X_*$ occurring here contain only
nonopen orbits. On a point orbit both
contractions are the identity. On every other
such orbit, (C.52.1) and C.50 use the same
actual intrinsic Lie image and the same
based target-cover path. Exponential naturality
therefore gives equality, including the
intermediate real circle angles described above.

Distinct open-orbit closures meet only in the
nonopen locus. The present construction assigns
there the same rule (C.52.1), so their contractions
and inverses agree. Conjugation by $h\in H$
identifies the actual target images and their
based paths, preserves metabelianness, and
intertwines both branches of (C.52.1).
Choose the intrinsic construction on one
representative of each $H$-orbit of components
and transport it by $h$. Commutation with the
full stabilizer makes this choice independent
of the transporting element, exactly as in
(C.51.15).

Finite closed-set pasting now gives the
homeomorphism families and inverse families
on (C.52.3). A general $h$ intertwines this
union with the union indexed by $h\mathcal I$;
when it preserves the union, it commutes with
the displayed contraction. The family on
$X_*$ is exactly the prescribed one throughout.
The diagonal proof concerns a single fixed
closure. It supplies no uniform control over
infinitely many different compact fundamental
sets, and no compact fundamental set for the
remaining noncocompact quotients. These are
still retained parts of C.12.
\end{proof}
